% Section 15: competitor analysis summary.
\section{Competitor analysis}\label{sec:competitors}

Before deciding what to build next, we installed and ran the established solvers on the development
machine and read their logs (\file{docs/COMPETITOR\_ANALYSIS.md}, driver \file{bench/comparators.py}).
The field covers HiGHS (dual simplex, IPX interior point, two PDLP variants, MIP, QP), SCIP with SoPlex,
Google OR-Tools (GLOP, PDLP, CLP, CBC, SCIP, CP-SAT), CBC, GLPK, Clarabel, OSQP and NVIDIA cuOpt on the
GPU. CPLEX, Gurobi and Xpress are commercial and were not installed. Each (instance, solver) pair ran in
its own process at default settings on small sets of \nCompLpInst{} Netlib LPs, \nCompMipInst{} MIPLIB~3
instances and \nCompQpInst{} Maros--Mészáros QPs. \Cref{tab:comp-lp,tab:comp-mip,tab:comp-qp} add
\nirnay{}'s own results on the same instances, taken from its sweeps.

\begin{table}[ht]
\centering\small
\caption[Comparator solvers on the Netlib LP set]{Comparator solvers on the LP set (\CompLpInstances).
Solved: objective within $10^{-6}$ of the reference ($10^{-3}$ for the first-order PDLP and OSQP runs,
whose default tolerances are loose). Times: shifted geometric mean over the solved instances and the
largest time. From \file{results/comparators\_lp.csv} and \nirnay{}'s sweeps; sorted by solved count,
then time.}
\label{tab:comp-lp}
\begin{tabular}{@{}l l r r r@{}}
\toprule
Solver & Platform & Solved & sgm (s) & max (s)\\
\midrule
\rows{data/comp_complp.tex}
\bottomrule
\end{tabular}
\end{table}

\begin{table}[ht]
\centering\small
\begin{minipage}[t]{0.52\linewidth}
\centering\footnotesize\setlength{\tabcolsep}{3pt}
\captionof{table}[Comparator solvers on the MIPLIB 3 set]{MIP set (\CompMipInstances); solved: optimal within $10^{-4}$. $^\ast$: GPU.}
\label{tab:comp-mip}
\begin{tabular}{@{}l r r r@{}}
\toprule
Solver & Solved & sgm (s) & max (s)\\
\midrule
\rows{data/comp_compmip_short.tex}
\bottomrule
\end{tabular}
\end{minipage}\hfill
\begin{minipage}[t]{0.45\linewidth}
\centering\footnotesize\setlength{\tabcolsep}{3pt}
\captionof{table}[Comparator solvers on the Maros--Mészáros set]{QP set (\CompQpInstances); solved: within $10^{-6}$ ($10^{-3}$ for OSQP).}
\label{tab:comp-qp}
\begin{tabular}{@{}l r r r@{}}
\toprule
Solver & Solved & sgm (s) & max (s)\\
\midrule
\rows{data/comp_compqp_short.tex}
\bottomrule
\end{tabular}
\end{minipage}
\end{table}

\subsection{What the comparison shows}
\begin{itemize}
  \item \textbf{Every winner is a dual simplex first.} HiGHS, SoPlex inside SCIP, GLOP and cuOpt's
        concurrent mode all rely on a dual simplex for small and medium LPs and for node LPs. The
        interior point is the alternative for large LPs, and PDLP a third engine for very large LPs at
        low accuracy. cuOpt races all three and takes the first answer, and on the small Netlib LPs its
        own concurrent mode always picked the CPU dual simplex. \nirnay{} has the same three engines;
        running them concurrently is on the roadmap.
  \item \textbf{PDLP failures are algorithmic.} Every PDLP implementation tested failed on the same
        badly scaled Netlib problems (\code{pilot4}, \code{greenbea}), as ours does. On large,
        structured LPs the picture reverses: GPU PDLP beats every factorisation-based method on the
        laptop (\cref{sec:pdlp-results}).
  \item \textbf{Presolve matters more than anything else.} A census of HiGHS's presolve on Netlib and
        MIPLIB~3 shows that singleton rows, doubleton equations, the aggregator (free-column
        substitution), forcing rows, and fixed and dominated columns do most of the work. It is also what
        makes an interior point robust: Clarabel and GLPK's interior point fail on \code{greenbea}
        without it, and HiGHS's IPX succeeds with it.
  \item \textbf{MIP gaps close at the root.} Gomory mixed-integer and (c-)MIR cuts are the common core
        of SCIP, CBC and cuOpt, with knapsack cover, zero-half, clique and implied-bound cuts next. First
        incumbents come from cheap heuristics run before or right after the first LP (HiGHS's
        feasibility jump, SCIP's trivial, locks and shifting heuristics).
  \item \textbf{References make mistakes too.} The study found a wrong ``optimal'' from one solver's
        MIP on \code{p0201}, a ``Solved'' status at a $10^{-3}$ objective error from another on
        \code{greenbea}, false infeasibility verdicts from GLPK's interior point, and a time-limit
        defect in one PDLP. \nirnay{}'s harness therefore scores against a reference in a separate
        process and records the reference's own status (\cref{sec:verification}).
\end{itemize}

\subsection{Licences and cost}
All the open-source comparators are free: HiGHS (MIT), SCIP and OR-Tools, Clarabel, OSQP and cuOpt
(Apache-2.0), CBC (EPL-2.0), and GLPK (GPL-3.0, whose copyleft prevents linking into a non-GPL product).
CPLEX, Gurobi and Xpress are proprietary: their free editions are limited in model size, and
production licences are sold per user or per deployment, with prices quoted on request. \nirnay{} is
intended to be released under Apache-2.0, like most of the open-source field, and because it depends
on no solver library, no licence restriction flows into it.
